% Part: intuitionistic-logic
% Chapter: tableaux
% Section: rules

\documentclass[../../../include/open-logic-section]{subfiles}

\begin{document}

\olfileid{int}{tab}{rul}

\olsection{అంతఃప్రజ్ఞావాద తర్క నియమాలు}

$\land$, $\lor$ సంధాయకాల నియమాలు సాధారణ
ప్రతిజ్ఞావాక్య చిహ్నిత \tetoken{టాబ్లోల}{!!{tableau}s}
నియమాలే; పూర్వసూచికలు మాత్రమే అదనంగా
ఉంటాయి. ప్రతి సందర్భంలో, చిహ్నిత
\tetoken{సూత్రం}{!!{formula}}
$\sFmla{S}{!A}[\sigma]$కు వర్తింపజేసిన నియమం,
అదే~$\sigma$ పూర్వసూచిక గల కొత్త
\tetoken{సూత్రాలను}{!!{formula}s}
ఇస్తుంది. ఇది అంతఃప్రజ్ఞాత్మకంగా స్పష్టం:
ఉదాహరణకు $!A
\land !B$ అనేది~$\sigma$ పేరుగల లోకంలో
సత్యమైతే, $!A$, $!B$ రెండూ~$\sigma$
వద్దే సత్యం (వేరే లోకంలో అని కాదు).
$\land$, $\lor$ నియమాలను
\olref{tab:prop-rules}లో సమీకరించాం.

\begin{table}
  \[\def\arraystretch{3}\begin{array}{|c|c|}
    \hline
    \AxiomC{\sFmla{\True}{!A \land !B}[\sigma]}
    \RightLabel{\TRule{\True}{\land}}
    \UnaryInfC{\sFmla{\True}{!A}[\sigma]}
    \noLine
    \UnaryInfC{\sFmla{\True}{!B}[\sigma]}
    \DisplayProof
    &
    \AxiomC{\sFmla{\False}{!A \land !B}[\sigma]}
    \RightLabel{\TRule{\False}{\land}}
    \UnaryInfC{$\sFmla{\False}{!A}[\sigma] \quad \mid \quad
      \sFmla{\False}{!B}[\sigma]$}
    \DisplayProof
    \\[2ex]
    \hline
    \AxiomC{\sFmla{\True}{!A \lor !B}[\sigma]}
    \RightLabel{\TRule{\True}{\lor}}
    \UnaryInfC{$\sFmla{\True}{!A}[\sigma] \quad \mid \quad
      \sFmla{\True}{!B}[\sigma]$}
    \DisplayProof
    &
    \AxiomC{\sFmla{\False}{!A \lor !B}[\sigma]}
    \RightLabel{\TRule{\False}{\lor}}
    \UnaryInfC{\sFmla{\False}{!A}[\sigma]}
    \noLine
    \UnaryInfC{\sFmla{\False}{!B}[\sigma]}
    \DisplayProof
    \\[2ex]
    \hline
  \end{array}\]
  \caption{$\land$, $\lor$లకు పూర్వసూచిక గల \tetoken{టాబ్లో}{!!{tableau}} నియమాలు}
  \ollabel{tab:prop-rules}
\end{table}

సంవృతం అయ్యే షరతు సాధారణ
\tetoken{టాబ్లోల}{!!{tableau}s}
షరతును పోలి ఉంటుంది; అయితే
\tetoken{సూత్రాలు}{!!{formula}s}
మాత్రమే కాక పూర్వసూచికలూ సరిపోవాలి.
నిజానికి కొంచెం ఉదారంగా ఉండవచ్చు:
అంతఃప్రజ్ఞావాద నమూనాల్లో
\tetoken{సూత్రాలు}{!!{formula}s}
ఒకసారి సత్యమైతే సత్యంగానే ఉంటాయి.
అందువల్ల~$\sigma$ వద్ద $!A$
సత్యమై, ప్రాప్యమైన ఏ
పూర్వసూచిక~$\sigma.{*}$ వద్దైనా
అసత్యమవడం అసాధ్యం. కాబట్టి ఒక శాఖలో
ఈ రెండూ ఉంటే అది సంవృతం:
\[
\sFmla{\True}{!A}[\sigma] \quad\text{and}\quad \sFmla{\False}{!A}[\sigma.{*}]
\]
ఏదో పూర్వసూచిక $\sigma$, \tetoken{సూత్రం}{!!{formula}}~$!A$
కోసం. చిహ్నాలు విరుద్ధ క్రమంలో ఉంటే,
అంటే కిందివి ఉంటే,
\[
\sFmla{\False}{!A}[\sigma] \quad\text{and}\quad \sFmla{\True}{!A}[\sigma.{*}]
\]
ఆ శాఖ సంవృతమయ్యేది $*$ శూన్య
క్రమమైనప్పుడే అని గమనించండి.

అదనంగా, శాఖలో~$\sFmla{\True}{\bot}[\sigma]$
ఉన్నా అది సంవృతం.

పరికల్పనలను అమర్చే నియమాలూ
సాధారణ \tetoken{టాబ్లోల}{!!{tableau}s}
నియమాలే; అయితే పరికల్పనలకు
ఎల్లప్పుడూ~$1$ పూర్వసూచికను వాడుతాం.
(అన్ని పరికల్పనలకు ఒకే పూర్వసూచికను
వాడితే, అది ఏదైనా పరవాలేదు.)
అందువల్ల, ఉదాహరణకు,
\[
!B_1, \dots, !B_n \Proves !A
\]
అని చెప్పేది, కింది పరికల్పనలకు
సంవృత టాబ్లో ఉన్నప్పుడే:
\[
\sFmla{\True}{!B_1}[1], \dots, \sFmla{\True}{!B_n}[1],
\sFmla{\False}{!A}[1].
\]

సోపాధికం~$\lif$ నియమాలు సాంప్రదాయిక,
మోడల్ తర్క కేసులకు భిన్నం.
$\TRule{\True}{\lif}$ నియమం
$\sFmla{\True}{!A \lif !B}[\sigma]$
గల శాఖను రెండు వేర్వేరు శాఖలపై
$\sFmla{\False}{!A}[\sigma.{*}]$,
$\sFmla{\True}{!B}[\sigma.{*}]$
చేర్చి విస్తరిస్తుంది.
\sourcecorrection{OLTEINTTABRULE-001}{మూల గద్యంలో సత్య సోపాధిక నియమం రెండు ఫలితాల సత్య-చిహ్నాలను పట్టికకు విరుద్ధంగా ఇచ్చింది; నియమపు రెండు వాదనల క్రమాన్నీ తారుమారు చేసింది. గద్యాన్ని దిగువ పట్టికలోని అసత్య పూర్వపక్షం లేదా సత్య ఉత్తరపక్షం అనే సరైన శాఖలతో, అదే నియమ సంకేత క్రమంతో సరిపోల్చాం.}
ఈ నియమాన్ని వర్తింపజేసే శాఖలో
\emph{ఇప్పటికే} ఉన్న పూర్వసూచిక~$\sigma.{*}$
కోసమే వర్తింపజేయవచ్చు. అలాంటి
పూర్వసూచికను ఆ శాఖపై ``వాడినది''
అందాం. ($\sigma.{*}$లో $\sigma$
కూడా ఉంటుంది కాబట్టి,
$\sFmla{\False}{!A}[\sigma]$,
$\sFmla{\True}{!B}[\sigma]$ అనే
పూర్వసూచిక గల చిహ్నిత సూత్రాలను
వేర్వేరు శాఖల్లో చేర్చి ఈ నియమాన్ని
ఎప్పుడూ వర్తింపజేయవచ్చు.)

$\TRule{\False}{\lif}$ నియమం
$\sFmla{\False}{!A \lif !B}[\sigma]$
గల శాఖపై
$\sFmla{\True}{!A}[\sigma.n]$,
$\sFmla{\False}{!B}[\sigma.n]$
రెండింటినీ అదే శాఖకు చేరుస్తుంది;
ఇక్కడ $\sigma.n$ ఆ శాఖకు కొత్త
పూర్వసూచిక.

$\lnot$ నియమాలను కూడా
ఇదే విధంగా నిర్వచిస్తాం
($\lnot !A$ను $!A \lif \lfalse$గా
నిర్వచించడాన్ని ఉపయోగించి).

నియమాలు \olref{tab:rules-lif-lnot}లో ఉన్నాయి.

\begin{table}
  \[\def\arraystretch{3}\begin{array}{|c|c|}
    \hline
    \AxiomC{\sFmla{\True}{\lnot !A}[\sigma]}
    \RightLabel{\TRule{\True}{\lnot}}
    \UnaryInfC{$\sFmla{\False}{!A}[\sigma.{*}]$}
    \DisplayProof
    &
    \AxiomC{\sFmla{\False}{\lnot !A}[\sigma]}
    \RightLabel{\TRule{\False}{\lnot}}
    \UnaryInfC{\sFmla{\True}{!A}[\sigma.n]}
    \DisplayProof
    \\[1ex]
    \text{$\sigma.{*}$ వాడినది} & \text{$\sigma.n$ కొత్తది}\\
    \hline
    \AxiomC{\sFmla{\True}{!A \lif !B}[\sigma]}
    \RightLabel{\TRule{\True}{\lif}}
    \UnaryInfC{$\sFmla{\False}{!A}[\sigma.{*}] \quad \mid \quad
      \sFmla{\True}{!B}[\sigma.{*}]$}
    \DisplayProof
    &
    \AxiomC{\sFmla{\False}{!A \lif !B}[\sigma]}
    \RightLabel{\TRule{\False}{\lif}}
    \UnaryInfC{\sFmla{\True}{!A}[\sigma.n]}
    \noLine
    \UnaryInfC{\sFmla{\False}{!B}[\sigma.n]}
    \DisplayProof
    \\[1ex]
    \text{$\sigma.{*}$ వాడినది} & \text{$\sigma.n$ కొత్తది}\\
    \hline
  \end{array}\]
  \caption{$\lnot$, $\lif$లకు పూర్వసూచిక గల \tetoken{టాబ్లో}{!!{tableau}} నియమాలు}
  \ollabel{tab:rules-lif-lnot}
\end{table}

\end{document}
